\section{Scope, comparison methods, and remaining questions}
\label{context:scope}

\paragraph{Relativistic matter and the degrees of freedom.}
Lieb--Yau~\cite{LY1988} treat relativistic particle kinetic energy
with Coulomb interactions and point nuclei, without a pair-creating
Dirac field or dynamical photons. The quantized-field theorem of
Lieb--Loss~\cite{LL2002} is a no-pair theory: its electron space
uses the positive spectral subspace of a field-dependent Dirac
operator. It admits fixed point nuclei and quantized photons but
excludes positrons. At physical $\alpha$ their stated range includes
$Z\le42$. Lieb's review~\cite{Lieb2002} explains these distinctions.
The magnetic critical-charge result of
Frank--Lieb--Seiringer~\cite{FLS2007} likewise does not add a
pair-creating Dirac sea to this argument.

J\"akel--Wreszinski~\cite{JW2018} consider finite-volume,
ultraviolet-regularized Coulomb-gauge QED with the fermion and photon
fields. Their finite-regulator vacuum-energy shift is allowed to
depend on volume. It is not the point-nucleus comparison uniform
in $K$ and volume proved in the present parameter ranges.
No no-pair or fixed-volume lower bound is used here as a substitute
for that comparison.

\paragraph{Relation to Gaussian domination in lattice models.}
Miyao's theorem~\cite{Miyao} concerns a half-filled Hubbard model on a
bipartite lattice, with two spin species, positive-type density interactions,
and quantized photons entering hopping amplitudes through Peierls phases.
Its finite-volume Theorem~3.4 proves Gaussian domination for arbitrary real
source shifts.  The proof uses a particle--hole transformation in one spin
sector and conjugate fermionic factors.  The thermodynamic susceptibility
theorem additionally assumes a summable interaction and yields
$\chi_\beta(p)\leq\widehat U(p)^{-1}$ where $\widehat U(p)>0$.
Unscreened Coulomb interactions do not satisfy that summability assumption.
The spin-mixing Dirac kinetic energy, projected charge operators, Wick
corrections, and fixed-charge projection in the present model require a
separate argument.  Thus the determinant lemma proved here is the actual
input; the lattice theorem provides methodological context rather than an
application covering this Hamiltonian.

Lieb--Siedentop~\cite{LS} already distinguish a positive density-product
interaction from its normally ordered form and analyze the resulting
quadratic term using a self-consistent choice of normal ordering.  Here
the reference is the prescribed free sea, and the density-product
rewriting retains its exact quadratic and scalar corrections.  Their
renormalization construction is not being substituted for this fixed
ordering.  The matter result of Bach--Barbaroux--Helffer--Siedentop~\cite{BBHS}
concerns a Hartree--Fock functional with a classical magnetic field;
it does not supply a lower bound over all correlated fermionic states
and quantized photon states.

\paragraph{Why projected densities need not commute.}
For scalar functions $f,g$, the matrices $P_{\Lambda_f}fP_{\Lambda_f}$ and
$P_{\Lambda_f}gP_{\Lambda_f}$ need not commute.  Accordingly, a Gaussian
identity for a simultaneous exponential of several noncommuting charges
would be unjustified.  The valid construction first writes the positive
Coulomb quadratic form as $\frac12\sum_a u_a R_a^2$, with Hermitian $R_a$
and $u_a>0$, and uses the identity
\begin{equation}
 e^{-\delta u_a(R_a-h_a)^2/2}
 =\int_{\mathbb R}\frac{e^{-z^2/2}\,dz}{\sqrt{2\pi}}
   e^{i\sqrt{\delta u_a}zR_a}
   e^{-i\sqrt{\delta u_a}zh_a}
 \label{context:hs}
\end{equation}
for each component separately.  The spectral theorem proves this
single-operator identity.  The finite-dimensional Trotter formula then
retains the order of all factors.  Positivity of their centered fermionic
trace, including the charge-projection phase, is supplied by the
positive-mass determinant lemma.  No commutation between different $R_a$
is assumed.  This argument does not establish an operator-valued
Graf--Schenker localization inequality.  The subsequent heat-kernel
comparison concerns a prescribed classical nuclear measure and the
bounded projected charge density.

\paragraph{Finite sources versus linear response.}
An alternative response argument would need control along the whole
source path.  For a fixed charge sector, put
$f(t)=\inf\sigma(H_0+tW)$, $0\leq t\leq1$.  This function is concave.
With $\mu=-f''$ interpreted as its nonnegative distributional second
derivative measure,
\begin{equation}
 f(1)=f(0)+f'_+(0)-\int_{(0,1)}(1-t)\,\mu(dt).
 \label{context:response-path}
\end{equation}
At a simple isolated ground state $\psi_t$, its ordinary density is
\[
 -f''(t)=2\big\langle (W-\langle W\rangle_t)\psi_t,
 (H_0+tW-f(t))^{-1}_{\perp}
 (W-\langle W\rangle_t)\psi_t\big\rangle.
\]
Level crossings contribute atoms to $\mu$.  The example
$H_0=\operatorname{diag}(0,\Delta)$,
$W=\operatorname{diag}(0,-B)$, with $B>\Delta>0$, has zero response at
$t=0$ but $f(1)=\Delta-B$.  Consequently a susceptibility bound only at
the unsourced vacuum cannot replace finite-source domination.  Away
from zero source, translation symmetry is also lost, so a diagonal
Fourier bound need not control the full response matrix.  The canonical
Gaussian-domination argument in this paper avoids these additional
hypotheses within its stated positive-mass regime.

\paragraph{The direction of a screened comparison.}
For $k\ne0$, write
\[
 \frac{4\pi}{|k|^2}=\frac{4\pi}{|k|^2+\mu^2}+w_\mu(k),\qquad
 w_\mu(k)=\frac{4\pi\mu^2}{|k|^2(|k|^2+\mu^2)}.
\]
The density-product remainder
$R_\mu=\frac{\alpha}{2V}\sum_{k\ne0}w_\mu(k)\rho(-k)\rho(k)$
is positive.  Its fully Wick-ordered counterpart need not be positive;
the associated quadratic and scalar corrections must remain in the
auxiliary screened Hamiltonian $\widetilde H_Y$ when writing the exact
identity $H_C=\widetilde H_Y+R_\mu$.  If $e(t)$ is the ground energy of
$\widetilde H_Y+tR_\mu$, the variational principle gives, for endpoint
ground states,
\begin{equation}
 \langle R_\mu\rangle_C\leq e(1)-e(0)
 \leq\langle R_\mu\rangle_Y.
 \label{context:screen-direction}
\end{equation}
Thus an upper bound on the Coulomb-vacuum expectation is not an upper
bound on the vacuum-energy difference that must be subtracted.  One
needs a screened-vacuum bound or an appropriate bound along the
interpolation.

The required estimate also concerns an equal-time structure factor.
For a translation-invariant simple ground state, use the conventions
\[
 S(k)=V^{-1}\langle\rho(-k)\rho(k)\rangle,
 \qquad
 \chi(k)=V^{-1}\langle\rho(-k)(H-E)^{-1}_{\perp}\rho(k)\rangle.
\]
Then $\langle R_\mu\rangle=\frac\alpha2\sum_k w_\mu(k)S(k)$.
The hypothesis $S(k)\leq s|k|^2$, on the finite band of charge modes,
gives $\langle R_\mu\rangle\leq C\alpha s\mu^2\Lambda_f V$ by the
lattice bound for $\sum_{k\ne0}|k|^{-2}$.
A bound $\chi(k)\leq s|k|^2$ alone does not imply this hypothesis.
Indeed, with
$F(k)=V^{-1}\langle\rho(-k)(H-E)\rho(k)\rangle$,
Cauchy--Schwarz gives $S(k)^2\leq\chi(k)F(k)$.
An additional estimate $F(k)\leq f|k|^2$ would suffice; it is not
silently available for the projected interacting theory.  A lower
spectral gap bounds $\chi$ in terms of $S$, in the opposite direction
from the inference just ruled out.

\paragraph{Vacuum normalization and source-species constants.}
A common scalar vacuum subtraction cancels from every vacuum-relative
energy.  In contrast, adding
$\sum_{i=1}^K\varepsilon_{Z_i}(\Lambda_f)$ changes the source energy
reference.  Such constants can be meaningful only if each nuclear
species has an explicitly stated, independently fixed calibration or
mass-renormalization condition.  They must be independent of the state,
positions, volume, and number of nuclei.  Their introduction is an
additional model choice for prescribed infinitely massive sources;
it is distinct from removing the infinite classical self-energy of
each point nucleus.  Defining the constants from the exact energy in
each finite box would generally introduce forbidden volume dependence.
Even a valid one-nucleus calibration would not establish a bound on
many-nucleus cross terms. No additional nuclear mass-calibration
subtraction is used in the momentum-cutoff theorems. The separate
lattice theorem in Section~\ref{lat-section} explicitly prescribes
individual electrostatic diagonals, using either the periodic Green
function or its volume-independent infinite-lattice counterpart.

\paragraph{Quantized photons and the energy reference.}
Section~\ref{pe-section} extends the comparison to the original
finite-mode quantized transverse field.  In real standing-wave
coordinates its normally ordered free Hamiltonian is
\[
 H_\gamma=\frac12\sum_r
 (-\partial_{X_r}^2+\omega_r^2X_r^2-\omega_r),\qquad
 \omega_r=|k_r|>0.
\]
The free-sea current has zero trace, so the current coupling adds only
Hermitian vector coefficients to the one-particle matrix.  The positive
mass coefficient is unchanged.  Componentwise Gaussian integration,
positive Mehler kernels, trace-norm Trotter convergence for a bounded
coordinate truncation, and its form-controlled removal prove the
canonical comparison with the \emph{interacting photon vacuum}.
The oscillator zero-point subtraction is the same on both sides.
No positive bare photon mass or fixed infrared decoupling is introduced.

Every mode with $0<|k|\leq\Lambda_\gamma$ is retained; hence arbitrarily
soft modes occur as $L\to\infty$.  The harmonic transverse $k=0$ sector
is explicitly fixed, rather than eliminated by Coulomb gauge alone.
The many-nucleus restriction $L\geq\Lambda_f^{-1}$ concerns the
periodic source estimates, not a photon infrared cutoff; the
one-nucleus estimate holds for every admissible box.  The fermion
cutoff in canonical momentum is not gauge covariant.  This remains
a specified regulated Hamiltonian, not a construction of renormalized,
gauge-invariant continuum QED.

In particular, a bound $E_\nu-E^{\rm vac}\geq-CK$ alone places no
restriction on the absolute vacuum baseline: a common, source-independent
scalar added to both Hamiltonians cancels in the difference.
The momentum-cutoff relative theorem neither proves absolute vacuum
extensivity nor establishes a thermodynamic limit with its given
normalization. The compact lattice vacuum limit in
Section~\ref{therm4-section} uses additional local cutting estimates.

\paragraph{What the result does not settle.}
The estimate now includes all correlated fermionic and photon states
at the stated cutoffs.  The sufficient positivity condition on the
effective mass is a condition on the bare parameters; it does not
control removal of the fermion cutoff at fixed physical mass and
charge.  BDF vacuum existence~\cite{HLS}
and charge renormalization in the reduced BDF model~\cite{GLS} are
results for different nonlinear mean-field functionals, not a
renormalization construction for this many-body Hamiltonian.
The Pauli--Villars construction with classical electromagnetic
fields~\cite{GHLS} likewise does not provide a theorem with quantized
transverse photons.

A low-energy Pauli--Fierz or Coulomb approximation would require control
of the relevant pair-production threshold and effective-operator
remainders, with errors extensive in system size.  Stability of an
effective Hamiltonian alone does not transfer to the original theory.
Finally, a lower bound is not an existence proof for a thermodynamic
limit, atomic binding, molecular binding, or crystalline equilibrium.
Those questions require separate compactness, binding, and structural
arguments. The lattice vacuum and open-source limits proved here have
the explicit boundary conditions of Section~\ref{therm4-section};
they do not assert the remaining periodic-source comparison.

\paragraph{The scope of Fry's strong-field calculation.}
Fry~\cite{Fry} studies a renormalized Euclidean fermion determinant
in a specified class of static magnetic fields. Section IV of the
revised manuscript explicitly leaves two assumptions to be established:
the continuum limit of the lattice diamagnetic inequality and the
asymptotic Pauli heat-kernel expansion. Its regularity and field-class
conditions must also be retained. We therefore do not use that
calculation as an unconditional instability theorem for the present
Hamiltonian or as a no-go theorem for every calibrated,
cutoff-uniform source-relative bound. The concrete limitation of
our estimate is its explicit $a^{-1}$ or $\Lambda_f$ source cost.
